%! AP Statistics — Unit 5, Lesson 5.2 — Group Activity
\documentclass[10pt]{article}
\usepackage{apstats-article}
\usepackage{apstats-boxes}
\newcommand{\TallMath}[1]{$\displaystyle #1\rule[-1.4em]{0pt}{3.2em}$}

\begin{document}

\pageheader{Unit 5, Lesson 5.2}{Group Activity}
\namepartnerperiod

\vspace{0.15in}

\begin{tcolorbox}[colback=redbg, colframe=black!40, boxrule=0.6pt, arc=2mm,
    left=3mm, right=3mm, top=2mm, bottom=2mm, title=\textbf{Tier R --- Remediate}]
\textbf{Context: Baby weights.} Newborn birth weights in the U.S. are approximately
normal with mean $\mu = 7.5$ lb and standard deviation $\sigma = 1.1$ lb.

\begin{enumerate}
  \item Use the 68--95--99.7 rule to find the range of birth weights for the middle
        95\% of newborns. Show your calculation.
        \writelines{2}
  \item A newborn weighs 6.0 lb. Compute the z-score and determine whether this weight
        is unusually low. Show work.
        \writelines{3}
  \item Using the z-score from part (b), the probability of a birth weight \emph{below}
        6.0 lb is approximately 0.086. Interpret this probability in context.
        \writelines{2}
\end{enumerate}
\end{tcolorbox}

\vspace{0.1in}

\begin{tcolorbox}[colback=redbg, colframe=black!40, boxrule=0.6pt, arc=2mm,
    left=3mm, right=3mm, top=2mm, bottom=2mm, title=\textbf{Tier A --- Approaching Proficiency}]
\textbf{Context: Standardized test.} Scores on a state exam are approximately normal
with $\mu = 72$ and $\sigma = 9$.

\begin{enumerate}
  \item Find $P(X > 85)$. Sketch, compute z-score, and use table/technology.
        \vspace{1.1in}
  \item Find $P(60 < X < 80)$. Show all steps.
        \vspace{1.1in}
  \item The top 10\% of scores earn a scholarship. What is the minimum score needed?
        Use inverse normal.
        \vspace{0.9in}
  \item Write a complete sentence interpreting your answer to part (c) in context.
        \writelines{2}
\end{enumerate}
\end{tcolorbox}

\vspace{0.1in}

\begin{tcolorbox}[colback=redbg, colframe=black!40, boxrule=0.6pt, arc=2mm,
    left=3mm, right=3mm, top=2mm, bottom=2mm, title=\textbf{Tier E --- Extension}]
\textbf{Context: Daily rainfall.} A weather model predicts that daily rainfall in a
city follows $N(1.2\text{ cm}, 0.4\text{ cm})$ during summer.

\begin{enumerate}
  \item Find the probability that a randomly selected summer day has rainfall between
        0.8 cm and 2.0 cm.
        \vspace{0.9in}
  \item Find the 5th and 95th percentiles. What do these values mean in context?
        \vspace{0.9in}
  \item The most extreme 2\% of rainfall days (highest and lowest) trigger weather
        alerts. Find the boundary values. (\textit{Hint:} split 2\% into two tails ---
        VAR-6.B.2d.)
        \vspace{0.9in}
  \item Is the normal model appropriate for daily rainfall? What features of the data
        would you check before applying this model?
        \writelines{3}
\end{enumerate}
\end{tcolorbox}

\end{document}
